
\chapter{Introduction}
\label{ch:intro}

Machine-learning (ML) models have reached state-of-the-art skill in weather
forecasting, rivaling traditional physics-based numerical models at a fraction
of their computational cost
\parencite{lam2023graphcast,bi2023pangu}.

A subfamily of these models are generative probabilistic weather models \parencite{price2023gencast,couairon2024archesweather}.
They are trained to approximate the conditional distribution of future atmospheric states.
This allows them to represent forecast uncertainty explicitly, alleviating the
smoothing effects often observed in deterministic approaches. Furthermore they enable
direct ensemble generation through the sampling, and achieve better probabilistic skill as
measured, for instance, by the Continuous Ranked Probability Score (CRPS)
\parencite{hersbach2000crps,price2023gencast}.

Beyond their use as probabilistic samplers, 
generative models can also be actively controlled during generation.
This is standard in the image generation domain, where applications such as text-to-image generation, editing,
and inpainting rely on \emph{guidance} to steer a pretrained
generative model toward a desired outcome
\parencite{dhariwal2021diffusion,ho2022classifierfree,nichol2021glide}.
Generative weather models, by contrast, are still predominantly used in their
unguided form, and their potential for guided generation remains
comparatively unexplored.

A small number of recent studies have investigated different forms of
weather control for generation of extreme events scenarios, including generation of tropical cyclones \parencite{brenowitz2025climate},
draughts
\parencite{ueyama2026guided}, and heatwaves
\parencite{whittaker2025heatwaves}.

This opens up a complementary use of generative weather models beyond
probabilistic forecasting.
While forecasting asks which future weather trajectories are likely given the
current state, \emph{scenario generation} asks what alternative weather
evolutions could look like under different targets while remaining meteorologically plausible.
This distinction is particularly relevant for the use case approached by the studies above: extreme weather storylines.
Extreme events account for a disproportionate share of weather-related impacts
\parencite{wmo2023atlas}, yet their rarity implies that both historical
observations and finite forecast ensembles provide a limited set of examples.
Yet, risk analysis requires considering plausible alternatives, 
including scenarios more severe than those previously observed, and once projected into the future, that live on the tails of the forecasting distribution
\parencite{thompson2017unseen,shepherd2018storylines}.

We are therefore interested in exploiting the ability of generative weather
models to produce coherent weather trajectories while deliberately steering
them toward extreme events.
% Rather than generating an isolated extreme state, the objective is to obtain an
% event that remains structured as weather throughout its evolution: coherent
% over time and space, across atmospheric pressure levels, and between the
% different meteorological variables.
% In other words, the imposed extreme should remain embedded within a
% multivariate atmospheric evolution that is consistent with the structure of
% realistic weather.
Transferring guidance methods from the image generation domain to this setting is nonetheless not
immediate.
Firstly, the conditioning objective has no ground truth, it has to be specified "by hand".
Secondly, the target is multidimensional, which makes the plausibility of outputs is considerably more structured to evaluate:
the generated trajectory must be coherent spatially, across pressure levels, variables have to covariate in a meaningful way, and the time evolution must also be realistic.

These considerations lead to the central research question of this thesis:

\begin{quote}
\emph{Can we steer generative probabilistic weather models to produce extreme
weather trajectories that satisfy a prescribed target while remaining
meteorologically plausible?}
\end{quote}

We tackle this question by identifying a gap in the current literature, 
and decompose the problem in two stages accordingly: specify the target and then guide.
Current approaches define a general target to be minimized or maximized, and play with different guidance strengths to obtain different scenarios.
We instead state that we should specify the target we want to achieve a priori, 
and derive the optimal guidance strength instead.

We investigate our research question using ArchesWeather
\parencite{couairon2024archesweather}, a medium-range probabilistic weather model based on
flow matching.
As experimental case study, we focus on regional surface-temperature
heatwaves and formulate the guidance objective in terms of a prescribed
temperature trajectory over a chosen region.
However, the method we develop is variable agnostic and could in principle be reformulated to generate different types of extreme weather storylines.
Furthermore, although extreme events provide the focus of this thesis, the underlying idea
is more general: guidance can be used to explore alternative weather
trajectories that deliberately depart from the base forecast toward any
suitably defined target.

In this thesis we make following main contributions:

\begin{enumerate}[noitemsep]
\item We formalize the notion of guidance strength schedule by unifying the current literature under this perspective.
\item We define an heuristic to specify the target scenario as a scalar trajectory of a
selected atmospheric variable averaged over a prescribed spatial region.
\item We develop a target-driven, gradient-based guidance method that derives the guidance strength to achieve the target under different schedules (how much should we introduce prior to make understand what we mean with this statement?). 
\item We identify the guidance strength schedule that achieves the most desirable results empirically and use this setting to construct different storylines.
\end{enumerate}


The remainder of the thesis is organized as follows.
\Cref{ch:background} introduces the technical background on flow matching and
guidance, reviews existing approaches to controlled weather generation, and
dissects ArchesWeather, before introducing the develped guidance method in \Cref{ch:specify,ch:guide}.
\Cref{ch:evaluation} presents the evaluation framework and 
\Cref{ch:experiments} presents the experimental results, which are discussed in
\Cref{ch:discussion}.
Finally, \Cref{ch:conclusion} summarizes the main findings and outlines
directions for future work.